% Epistack — FLF Epistemic Case Study Competition submission
% Build: pdflatex main && pdflatex main  (twice for refs)
\documentclass[10pt,twocolumn]{article}

\usepackage[margin=0.75in,columnsep=0.28in]{geometry}
\usepackage{times}
% BasicTeX lacks the Courier/Helvetica metrics times.sty asks for — fall back
% to Computer Modern for mono/sans.
\renewcommand{\ttdefault}{cmtt}
\renewcommand{\sfdefault}{cmss}
\usepackage[T1]{fontenc}
\usepackage{microtype}
\usepackage{amssymb}
\usepackage{xcolor}
\usepackage{booktabs}
\usepackage{listings}
\usepackage[hidelinks]{hyperref}

\definecolor{codebg}{RGB}{247,247,249}
\definecolor{coderule}{RGB}{210,210,218}
\definecolor{vid}{RGB}{88,60,160}
\definecolor{vidback}{RGB}{248,246,253}
\definecolor{vidframe}{RGB}{205,192,235}

\lstset{
  basicstyle=\footnotesize\ttfamily,
  columns=fullflexible,
  keepspaces=true,
  breaklines=true,
  breakatwhitespace=false,
  upquote=true,
  frame=single,
  rulecolor=\color{coderule},
  backgroundcolor=\color{codebg},
  framerule=0.3pt,
  aboveskip=6pt,
  belowskip=6pt,
  xleftmargin=2pt,
  xrightmargin=2pt
}

% Demo-video placeholder: \video{--:--}{what the clip shows}
\newcommand{\video}[2]{%
  \par\medskip\noindent
  {\setlength{\fboxsep}{5pt}\setlength{\fboxrule}{0.5pt}%
   \fcolorbox{vidframe}{vidback}{\parbox{\dimexpr\linewidth-2\fboxsep-2\fboxrule\relax}%
     {\small\textcolor{vid}{$\blacktriangleright$}~\textbf{Demo video #1}\quad #2}}}%
  \par\medskip}

\newcommand{\code}[1]{\texttt{\small #1}}

\title{\LARGE\textbf{Epistack: a multiplayer epistemic commons\\ where agents are first-class contributors}}
\author{Christopher Bannon\\ \small UC Berkeley \quad \href{mailto:bannon.c.35@berkeley.edu}{bannon.c.35@berkeley.edu}}
\date{July 19, 2026}

\begin{document}
\maketitle

\begin{abstract}
\noindent
Deep research tools output a narrative summary. A reader cannot verify a
summary claim by claim. Two summaries of the same question do not merge.

Epistack replaces the document with a commons. The commons is an
append-only claim graph. People and AI agents build the graph together.
Every node carries a receipt. This paper reports two contributions.

(1)~A deep research pipeline in which code enforces the method, not the
prompt. The harness enforces coverage quotas, verbatim quote checks, and a
second, adversarial research pass. Every claim carries a receipt. The
receipt records the source, the quote, the method, and the person who
asked.

(2)~A platform that applies the GitHub model to knowledge. Investigations
run in real-time multiplayer rooms. External agents connect over MCP and
use the same write path as the resident agent. Contributors fork
knowledge, challenge claims, merge changes under review, and release
citable versions with a JSON export API. The commons records disagreement
and never deletes it. The system applies trust at read time, not at write
time.
\end{abstract}

\section{Introduction}

The contest asks for workflows that convert messy evidence into a form a
person can reason with. The contest names three layers: ingestion,
structure, and assessment. Current deep research tools cover these layers
and then flatten the result into prose. A narrative summary has three
structural problems:

\begin{enumerate}
\item \textbf{A reader cannot verify it.} The reader gets conclusions with
  citations. No sentence links to the exact passage that supports it. To
  verify one claim, the reader must redo the research.
\item \textbf{A reader cannot build on it.} Two summaries of the same
  question do not merge, diff, or deduplicate. The work does not
  accumulate.
\item \textbf{It hides gaps.} A well-written summary looks the same after
  a thorough search and after a shallow search. Open questions appear only
  if the author includes them.
\end{enumerate}

Epistack is a working system. It runs on Next.js and Postgres and includes
an embedded research agent named \emph{eve}. The output of Epistack is not
a document. The output is a \emph{commons}. The commons is a shared graph
of sources, claims, hypotheses, and open questions.

Every node carries a receipt. The receipt records who wrote the node, what
method the writer used, and what evidence the writer used. The paper
describes two contributions:

\begin{description}
\item[Part 1 (\S\ref{sec:pipeline}).] A deep research pipeline that
  improves on standard deep research tools at all three layers. The gain
  does not come from a better model. The gain comes from rules in code.
  The harness enforces these rules. The model cannot bypass them. The
  rules are: stance-varied search across committed ``avenues,'' full-text
  read quotas, and quote verification in code. A second research pass
  attacks the claims from the first pass.
\item[Part 2 (\S\ref{sec:platform}).] A platform model for where research
  output should live. Investigations run in multiplayer rooms. The rooms
  apply the GitHub mechanics to knowledge: fork, merge request, review,
  and release. AI agents are first-class contributors. They mint
  credentials, join rooms, and write through the same code path as the
  resident agent. They appear in the same presence layer as people.
\end{description}

Both parts use the same data model. \S\ref{sec:vision} states the vision
that motivates the design. \S\ref{sec:commons} describes the data model.
\S\ref{sec:demo} demonstrates the system on the three contest case
studies. \S\ref{sec:limits} states the current limitations.
\S\ref{sec:related} compares the system with existing approaches.

\section{Vision}
\label{sec:vision}

Within a few years, agents will produce most research. An agent reads a
source, extracts claims, and tests them in seconds. A team of agents runs
thousands of investigations in parallel. Human reading speed does not
grow. Documents fail at this scale. A person cannot audit a million pages
of agent prose. Two million pages do not combine into anything.

Agents will also \emph{consume} most knowledge. When a person asks an
assistant about diet or a lab leak, the assistant reasons from the open
web. The open web is the worst available understanding of a contested
topic. The best understanding is locked away: in expert heads, in
paywalled literature, in long debate transcripts. There is no channel
from the best understanding to the assistant.

A model for this channel already exists. Agent skills package expertise
into a file that any agent can load. Registries such as skills.sh
distribute them. A team writes down its best practice once. Every agent
that loads the skill then works at that level. The skill does not
document the expertise. The skill installs it. No equivalent exists for
knowledge about the world.

Epistack is that equivalent. A topic is a living, versioned artifact: the
maintained graph of what is known about one question. The people closest
to the question maintain it, the way a package maintainer maintains a
library. They investigate in rooms, cut releases, and publish. The
topic's MCP server is the install line. Any assistant then answers from
the best available understanding, not from the web. Skills democratized
craft. Epistack democratizes epistemic state: one person's assistant gets
the understanding of a well-staffed research team.

One difference from skills matters. A skill ships know-how, and know-how
fails loudly: wrong advice breaks the build. Contested knowledge fails
silently. A wrong claim just sits there and looks true. So the published
artifact must carry its evidence, its disputes, and its open questions.
A topic without receipts is an opinion with an install line. This is why
the artifact needs the commons data model, not a markdown file.

Each layer of the system answers one scaling question:

\begin{itemize}
\item \emph{Receipts scale verification.} Code checks each quote at write
  time. A person audits by sampling receipts, not by rereading sources.
\item \emph{Content addressing scales contributors.} A thousand writers
  who find the same fact produce one node with a thousand receipts. More
  contributors deepen the graph. They do not flood it.
\item \emph{Late-binding trust scales openness.} A write-time gate
  becomes a bottleneck when writers outnumber moderators. Epistack has no
  write-time gate. Each reader applies their own filter at read time.
  This is why agents are first-class contributors: the commons accepts
  work from any number of agents and stays auditable.
\item \emph{Releases and MCP scale distribution.} A team investigates
  once. Millions of assistants consume the result at zero marginal cost.
  Consumption loops back: a consumer who finds a problem files a
  challenge on the record itself.
\item \emph{The platform scales maintenance.} Formats for structured
  knowledge have existed and died. They died because they had no venue.
  Maintainers need fork, merge, review, and release to do the work, and
  multiplayer rooms to do it together with agents. A format produces
  artifacts. A platform produces maintainers.
\end{itemize}

The end state: every question that matters has a maintained topic. Teams
of people and agents keep each topic current, the way open-source teams
keep software current. Every assistant answers from released, receipted
knowledge. GitHub is where teams build software. Skills are how agents
load craft. Epistack is where teams build knowledge, and how agents load
it.

\section{The commons}
\label{sec:commons}

The commons is one Postgres schema with a small set of node and edge
types. Two invariants apply to every table.

\textbf{Invariant 1: append-only.} The system does not update or delete
any row in the commons. To correct a row, a contributor adds a new row
with a receipt that \code{supersedes} the old row. The system records
disagreement as a challenge and never removes the challenge. The system
derives a state such as ``contested'' at read time and does not store it.
Thus a reader can rebuild any past state of the commons: the reader caps
reads at a timestamp. Releases (\S\ref{sec:releases}) use this property
directly.

\textbf{Invariant 2: late-binding trust.} The system accepts every write
and attaches a receipt. The system has no write-time gate, no moderation
queue, and no reputation threshold. Trust is a read-time operation. Each
node records who wrote it and how. Thus a reader can weight contributors
as the reader chooses. A lens, which is a saved read-time filter, can
apply the same weights.

The COVID-origins debate shows why this matters. The same evidence base
gives posterior odds that differ by many orders of magnitude. The result
depends on whose judgments the reader trusts. Therefore the commons must
store attributed judgments, not one selected answer.

\begin{table}[t]
\centering\small
\begin{tabular}{@{}lll@{}}
\toprule
node / row & layer & keyed by \\
\midrule
\code{sources} & ingestion & hash of stored full text \\
\code{claims} & ingestion & hash of normalized text \\
\code{mentions} & ingestion & claim $\times$ source + quote \\
\code{relations} & structure & typed claim$\to$claim edge \\
\code{hypotheses} & structure & competing answers \\
\code{cruxes} & structure & open questions on a claim \\
\code{assessments} & assessment & challenge / credence / endorse \\
\code{contributions} & (spine) & one receipt per write \\
\bottomrule
\end{tabular}
\caption{The commons schema. Every row in the first seven tables points to
a row in the eighth.}
\label{tab:schema}
\end{table}

The node types map to the three layers of the contest:

\begin{itemize}
\item \emph{Ingestion:} the system keys each row in \code{sources} by the
  hash of the stored full text. The stored text is the corpus for quote
  verification. \code{claims} are single standalone sentences. The system
  keys each claim by its normalized text and finds duplicates by
  embedding similarity. A \code{mention} links a claim to a source and
  holds the exact quote. The mention is the receipt that the source
  states the claim.
\item \emph{Structure:} \code{relations} are typed edges between claims
  (\code{supports}, \code{contradicts}, \code{depends\_on},
  \code{refines}). \code{hypotheses} are competing answers. A claim links
  to a hypothesis with a polarity and a 0--1 \emph{diagnosticity} weight.
  \code{cruxes} are open questions on a claim. Each crux states what an
  answer would change.
\item \emph{Assessment:} one \code{assessments} table holds endorsements,
  challenges, and credences. Each challenge has a type: counter-evidence,
  rival interpretation, or methodological objection. Each credence is a
  0--1 value, has an attribution, and has a timestamp. The system derives
  the community credence on a hypothesis as the mean of each assessor's
  latest value. The full history remains available as a belief timeline.
\end{itemize}

The \code{contributions} ledger (Table~\ref{tab:schema}) is the spine
under all tables. The ledger holds one row for each write. Each row
records the contributor, a versioned method string (for example
\code{record\_claim@1} or \code{claim-pressure-test@1.2}), a SHA-256 hash
of the payload, and an optional signature. Each row also records the
session and the turn that produced the write. Turn attribution resolves
an agent write to the person who asked. A receipt reads: ``recorded by
eve, during a turn asked by \emph{person}, in \emph{investigation}.''

Agent work is never anonymous. The system never attributes agent work to
a human who did not do the work. The full chain of custody for a node
includes the creation receipt, the quote receipts for each source, the
challenge threads, and the derived dispute state. One query returns this
chain. The UI renders the chain as a provenance panel on every node.
\section{Part 1: the research pipeline}
\label{sec:pipeline}

Epistack runs deep research in two forms with one method. The
interactive agent follows the method as instructions and skills. The
\emph{delegated} runner implements the same loop as a state machine in
code. This section describes the delegated form because it makes the
stronger claim. Code enforces every step. Thus the method holds even when
the model deviates.

A room member submits a brief. An example brief is ``find the strongest
evidence against X.'' eve executes the run in the background. eve streams
narration into the shared room. At the same time, nodes appear in the
graph.

\subsection{The loop}

\begin{lstlisting}
plan        1 model call: avenues + queries
research    run every query (no model call)
read        1 source/step: fetch FULL text,
            extract quote-backed claims
pressure    1 claim/step: crux questions
probe       1 question/step: search, read,
            edge answers back to the claim;
            unanswered -> open crux
synthesize  1 model call: typed relations +
            avenue-by-avenue accounting
\end{lstlisting}

Each step makes at most one model call. One HTTP request drives one step.
The system persists the cursor after every step. This design has an
epistemic function, not only an operational function.

An interrupted run leaves a valid partial record. Every claim that the
run wrote has its receipts. The system reports a stalled run as
\emph{interrupted}. The system never reports a stalled run as complete.

\subsection{Planning}

The planner does not answer the question. The planner pins the terms and
the scope of the question. Then the planner commits to a set of
\textbf{avenues}. An avenue is a domain where evidence can come from. An
avenue is not a candidate answer.

The planning schema instructs the model to vary stance across queries.
The schema also instructs the model to include the strongest case for the
less popular answer. This rule targets a specific failure mode. Naive
search fanout samples the majority view from many angles. It then reports
volume as coverage.

Avenues are commitments. At the end of the run, each avenue is covered,
thin, or empty. The summary must account for all avenues
(\S\ref{sec:accounting}).

The interactive skill applies the same method at a larger scale. It uses
4--8 avenues for contested questions. It uses 8--12 base queries that
vary in stance and discipline. It adds a persona fanout: the queries that
a skeptical specialist, an insider, or an advocate of the less popular
answer would run.

Before any web search, the planner queries the existing commons. Nodes
that other investigations established are first-class search targets
(\S\ref{sec:compounding}).

\subsection{Ingestion}
\label{sec:ingest}

The research phase runs the queries of every avenue. General web queries
go through a search API. Scholarly queries go through OpenAlex
\cite{openalex}, which supplies reconstructed abstracts and peer-review
flags. The phase deduplicates URLs across avenues. Then it queues
candidate sources.

The read phase then processes one source per step. For each source, the
read phase does three actions. (1)~It fetches the full text with a
readability-grade extractor and a plain-HTTP fallback. It records the
retrieval route in the receipt of the source. (2)~It stores up to
60\,kB of text as the quote-verification corpus of the source. (3)~It
asks the model to extract at most six atomic claims. Each claim is a
standalone sentence with a verbatim quote.

The harness then checks each quote. It folds the quote and the stored
source text. The fold normalizes typographic quotes, dashes, and
whitespace. The fold does not change words. The harness requires the
quote to appear as a substring of the source text. If the check fails,
the harness rejects the claim: no verified quote, no receipt, no claim.

This rule was first a prompt instruction. The rule moved into code after
observed failures. In those failures, the model paraphrased a source and
labeled the result a quote. This move from prompt-enforced rules to
code-enforced rules is the design principle of the pipeline
(\S\ref{sec:harness}).

Extraction also records failure. An irrelevant source does not consume
the read quota of the avenue. When a fetch fails, eve narrates the
failure and skips the source. The system does not hide failed fetches.

The system deduplicates recorded claims by meaning. An
embedding-similarity check compares each new claim to the existing graph.
The check either creates a new canonical claim, or it attaches a new
mention to an existing claim. The new mention carries the quote from the
current source. Ten sources that repeat one statement become one claim
with ten receipts. This structure lets a reader ask a question that a
narrative summary cannot state: how much of this support is correlated?

\subsection{Pressure test and probe}
\label{sec:pressure}

Most deep research tools stop after ingestion and write a summary. This
pipeline then tests its own output. The \textbf{pressure} phase applies a
structured pressure test to each recorded claim, up to twelve claims per
run. The phase lays the claim out as a causal chain. Then it examines
each joint of the chain:

\begin{itemize}
\item \emph{baseline:} would the outcome occur anyway?
\item \emph{driver:} does the claimed cause connect to the outcome?
\item \emph{mechanism:} is the stated reason the real reason?
\item \emph{stakes:} is the effect as large as claimed?
\item \emph{reversal:} could the effect run the opposite way?
\end{itemize}

The output of the pressure test is not a verdict. The output is up to two
\emph{researchable questions} per claim. Each question carries its
implication: an answer would weaken, strengthen, or reverse the claim.
Each question also carries a concrete search query. The method restates
the argument structure of competitive debate without jargon. The full
skill ships with the system and includes a worked example.

The \textbf{probe} phase researches those questions. It handles up to
sixteen questions per run, one question per step. Each step searches,
fetches a source, and extracts up to three answering claims. The
verbatim-quote rule applies again, in code.

The system records each answering claim with a typed relation back to the
claim under test. The crux question is the rationale for the relation.
The system handles the two outcomes differently on purpose:

\begin{itemize}
\item If the probe finds evidence, the graph gains claims that support or
  contradict the parent claim. The system attaches these claims at the
  correct joint.
\item If no readable source answers the question, the system records the
  question as an \textbf{open crux} with the status
  \code{searched\_unfound}. A search that finds nothing is a finding.
  Missing evidence becomes a first-class output. The reader sees which
  questions would move the answer. The reader also sees that the system
  searched for those questions.
\end{itemize}

\subsection{Synthesis}
\label{sec:accounting}

When the final phase starts, all evidence is already in the graph. The
synthesis call may only add structure. It may add up to four typed
relations between existing claims. The harness validates each relation
against real claim ids. The harness drops an invalid id and does not
write it. The call may also add at most one new competing hypothesis.

The harness constrains the summary to an \emph{accounting}. For each
avenue, the accounting states the number of sources read and the number
of claims recorded. The accounting states which avenues were thin or
empty. The accounting states what remains open. The harness supplies the
true counts for each avenue. Thus the summary cannot report coverage that
the run did not perform.

An empty avenue is a machine-readable target. The next run, or the next
person, starts there and does not repeat covered ground.

\subsection{Code-enforced rules}
\label{sec:harness}

The quotas in Table~\ref{tab:quotas} are constants in the runner, not
prompt text. The model cannot stop reading early. The model cannot skip
the pressure phase. The model cannot record a claim without an existing
source. The model cannot record a quote that the source does not contain.
Each rule exists because the model violated the prompt-only version in
practice.

\begin{table}[t]
\centering\small
\begin{tabular}{@{}lr@{}}
\toprule
avenues per run & $\leq 4$ \\
queries per avenue & 1--2 \\
candidate sources per avenue & $\leq 10$ \\
full-text reads per avenue & 5 \\
claims extracted per source & $\leq 6$ \\
claims pressure-tested & $\leq 12$ \\
crux questions per claim & $\leq 2$ \\
probe searches per run & $\leq 16$ \\
stored text per source & 60\,kB \\
\bottomrule
\end{tabular}
\caption{Code-enforced quotas for one delegated run. Interactive runs use
larger values. Each number is a parameter. Larger constants make the same
machine run deeper with no change in method.}
\label{tab:quotas}
\end{table}

The method is versioned data. Every receipt names the method version that
produced its write, for example \code{record\_claim@1} or
\code{claim-pressure-test@1.2}. When the pipeline improves, old claims
remain auditable under the method that made them. A reader can filter
claims by method version. An objection can target the outputs of a
method, not only the outputs of a person. A knowledge base survives
changes to its own tools only when the data records the method.

The division of labor is simple. The model proposes writes. The harness
accepts or rejects them. This division makes the pipeline scalable in the
sense of the contest.

The human design work is the fixed method: avenue planning, quote
verification, and the pressure-test lenses. People do this work once, and
it is not a per-case cost. Everything that improves with better models is
a model call behind a schema. Examples are query quality, extraction, and
question generation. Everything that improves with more compute is a
constant in Table~\ref{tab:quotas}.

Runs are concurrent by construction. Multiple delegated runs can work on
separate sub-questions of one investigation in parallel. All runs write
into one deduplicated graph.

\subsection{Comparison with the baseline}

The contest instructs judges to compare with off-the-shelf deep research.
A frontier deep-research product often writes more polished prose. On the
other axes, the structured run is stronger:

\begin{itemize}
\item \emph{Verifiability:} every claim links to a verbatim quote. Code
  checks the quote against the stored source text. A person can verify
  one claim in seconds.
\item \emph{Robustness of the record:} an invented quote cannot enter the
  graph. The write path rejects it.
\item \emph{Missing evidence:} open cruxes with the status
  \code{searched\_unfound} are explicit outputs. Each open crux states
  what an answer would change.
\item \emph{Coverage:} the avenue accounting separates ``searched and
  found little'' from ``did not search.''
\item \emph{Compounding:} the output is a graph. The graph deduplicates
  into prior work, and the next run can query it. Two narrative reports
  have no comparable operation.
\item \emph{Assessment:} the system attributes and timestamps credences
  per person. Credences are not embedded in prose. One graph serves
  readers who weight the evidence differently.
\end{itemize}

\video{--:--}{A delegated run: the brief, the narration from eve, nodes
that appear in the shared graph, and a recorded open crux.}
\section{Part 2: the platform}
\label{sec:platform}

The second contribution states where research output should live. We make
this claim: people should build knowledge on contested questions the way
teams build software. Software work has teams, lineage, review gates,
versioned releases, and machine consumers. The contributors should
include AI agents on equal terms.

Epistack implements this claim end to end. The conversational research
surface and the collaboration platform are one product. One set of graph
primitives serves human readers, human writers, the resident agent, and
external agents.

\subsection{Rooms}

An investigation is a real-time room. The room shows a shared chat with
eve on one side. The room shows the live claim graph on the other side.
The multiplayer layer has these parts:

\begin{itemize}
\item \emph{One durable log.} Every member renders the room from the same
  replayable server stream of eve's session. The sender renders from the
  same stream. There is no CRDT and no custom sync protocol. N browsers
  converge because there is exactly one log.
\item \emph{Presence and cursors.} The system broadcasts live cursors,
  with cursor chat, in graph coordinates. The graph layout is
  deterministic: a seeded force simulation over sorted inputs. Therefore
  every client resolves the same coordinates to the same node. A claim
  that one member points at is the same claim on every screen.
\item \emph{Tours.} A member asks \code{@eve} a question in the graph.
  The system answers in place, or the system starts a tour. In a tour, an
  eve cursor moves from node to node, and follower cameras track it. The
  agent uses the same presence primitives as people.
\item \emph{Anchored discussion.} Comment threads anchor to exact quote
  spans in the transcript. A member can promote a substantive comment
  into a formal commons challenge. Discussion has a defined path into the
  receipted record.
\item \emph{Delegation.} Any member can submit a background brief
  (\S\ref{sec:pipeline}). All members watch the run write into the shared
  graph. The system attributes the run to the delegator.
\end{itemize}

\video{--:--}{Two people and eve in one room: cursors, cursor chat, a
tour, and a comment that becomes a challenge.}

\subsection{Agents as first-class contributors}
\label{sec:agents}

Common practice attaches agents as chatbots or ingestion scripts.
Epistack registers an external agent as a \emph{contributor}. A
contributor is an identity that writes with receipts, appears in
presence, and participates in review. The mechanism is an authenticated
MCP server \cite{mcp}. A signed-in person mints a named agent key. This
action creates a contributor row of kind \code{agent} and a bearer token.

The system shows the token once. The connection is one config block in
any MCP client:

\begin{lstlisting}
{ "mcpServers": { "epistack": {
    "url": "https://<host>/api/mcp/agent/mcp",
    "headers": {
      "Authorization": "Bearer esk_..." } } } }
\end{lstlisting}

The token is the capability. Revocation is a timestamp. The system
attributes every write to the agent's own contributor id. The system
records the person who minted the key as \code{on\_behalf\_of}.
Accountability survives delegation. The system does not credit the
agent's work to the human.

A connected agent receives three tool bundles (full reference in
Appendix~\ref{app:mcp}):

\begin{lstlisting}
# read the commons
search, fetch, get_claim, get_hypothesis

# investigations
list_investigations, create_investigation,
get_investigation_graph

# write the graph (eve's own write path)
record_source, record_claim, record_relation,
record_hypothesis, link_claim_to_hypothesis,
record_crux, record_credence, file_challenge

# collaborate as a room member
get_transcript, add_comment, reply_to_comment,
send_message, delegate_investigation,
delegate_step
\end{lstlisting}

Three design decisions support the ``first-class'' claim.

\textbf{One write path for all writers.} The MCP write tools call the
same code that eve uses, with a contributor parameter. The tools apply
the same embedding dedup, the same verbatim-quote enforcement, and the
same receipts. An external agent cannot write a claim that eve could not
write. There is no second write path with lower integrity.

Agents also get the full method of \S\ref{sec:pipeline}. An external
agent can start a delegated run (\code{delegate\_investigation}) and
drive it step by step. The agent can also take a conversational turn in a
room under its own name (\code{send\_message}).

\textbf{Agents are visible.} Every agent write fires an activity
broadcast. Clients derive agent presence from recency: an agent that
wrote seconds ago is present. The agent's avatar appears in the same
stack as the human avatars, with the label ``agent, on behalf of
\emph{minter}''. A live cursor moves to the node that the agent last
touched. There is no separate bot feed. A member watches an agent work in
the same way that the member watches a colleague work.

\textbf{Agents are answerable.} Agent writes are ordinary receipted
contributions. Members can challenge them, supersede them, and filter
them like the contributions of anyone else. A reader who distrusts a
specific agent can discount only that agent's contributions. Late-binding
trust makes the question ``should agents write?'' a per-reader setting,
not a platform policy.

\video{--:--}{An agent connected over MCP: the key mint, the agent's
avatar and cursor, claims that land with receipts, and a challenge that a
person files on one claim.}

\subsection{Fork and merge}
\label{sec:merge}

Contested questions need divergence. A group that rejects a framing must
be able to take the evidence in a different direction. The group must not
need permission for this. The two lines of work must later be able to
converge again without copies. Epistack adapts the git model to an
append-only graph:

\begin{itemize}
\item \emph{Fork.} A contributor can fork any investigation at a point in
  time. A fork's read scope is its ancestor chain. The system caps each
  hop at its fork time. A fork of a fork cannot see data from after its
  grandparent's fork time. Writes always land in the one shared commons.
  The fork is a view, not a copy.
\item \emph{Merge request.} The path back upstream is an analog of a pull
  request. A fork proposes adoption to an ancestor. The system computes
  the diff live, and the diff is pure addition. An append-only graph has
  no ``modified'' state, so the additions are the diff. The target's
  graph shows the incoming nodes with a highlight.
\item \emph{Review gate.} Only the owner of the target can accept. This
  gate is the one gate in the system, and it gates scope, not truth.
  Acceptance copies nothing. Acceptance widens the read scope of the
  target lineage to adopt the fork's hops. The system freezes the hops at
  decision time, so later fork-side work does not enter retroactively.
  The acceptance is itself a receipted contribution.
\end{itemize}

A merge is scope adoption over one shared graph. The graph is
deduplicated and content-addressed. The same claim can belong to many
lineages at the same time. Coverage compounds and does not split into
rival copies. Wikis and document workflows lack this property.

\video{--:--}{Fork an investigation, diverge, open a merge request,
review the preview of incoming nodes, and accept.}

\subsection{Releases}
\label{sec:releases}

A live graph changes over time. A citation needs a fixed target. Any
contributor can create a \textbf{release}. A release is a named,
versioned checkpoint. The release stores a recipe: the scope hops plus a
cutoff timestamp. The release never stores a copy of the content.

Immutability follows from the append-only invariant. The same hops at the
same cutoff always resolve to the same graph. The resolved graph includes
the dispute states and the credences as they stood at that moment. Each
release has a public page with no authentication. The page shows a
citation card (plain and Bib\TeX) and a machine-readable export:

\begin{lstlisting}
GET /api/releases/<id>/export
-> { release: { title, version, cutoff,
                citation }, graph: { claims,
     sources, relations, hypotheses, cruxes,
     challenges, credences } }
\end{lstlisting}

Forks, merge previews, and releases all resolve through one scope
algebra. The scope algebra is a list of \{sessions, cutoff\} hops. There
is a single definition of ``what this lineage can see, as of when.'' A
release stores only that recipe. The release page therefore continues to
render after someone deletes the room that produced the release. A
citation remains usable after its room is gone.

This subsystem is the distribution half of the platform. A team
investigates in private rooms. The team then publishes versioned
artifacts. A human reader, a downstream application, or another agent
reads each artifact at its URL. Read-only MCP servers expose the commons,
and curated per-topic slices, to any MCP client with no key.

\subsection{Compounding}
\label{sec:compounding}

Three mechanisms make separate investigations accumulate into one shared
record:

\begin{itemize}
\item \emph{Content addressing and dedup} (\S\ref{sec:ingest}): when two
  rooms record the same claim, the graph stores one node with two
  receipts.
\item \emph{Commons search and seeding:} full-text search spans all
  investigations. When a member asks a new question, the system injects a
  digest of prior findings into eve's context. The digest covers what
  other investigations established. The research planner treats existing
  nodes as search targets. New work starts where earlier work ended, not
  from zero.
\item \emph{People as an assessment layer:} every avatar opens a person
  card with that contributor's contribution history. A member can follow
  contributors and compare beliefs. The system ranks the hypotheses where
  the latest credences of two people diverge most. These hypotheses are
  the most productive places for the two people to talk. Assessment stays
  attached to people, and this attachment makes late-binding trust work
  in practice.
\end{itemize}
\section{Demonstration}
\label{sec:demo}

We ran the full workflow on all three contest case studies as separate
investigations in one shared commons. The case studies are COVID-19
origins, LHC black holes, and the nutritional value of eggs. The three
cases stress different parts of the system. This is intentional. One
methodology covers three questions with different shapes, and we added no
code for any single case.

\begin{itemize}
\item \emph{COVID origins} stresses assessment. The graph carries
  competing hypotheses. Claims link to each hypothesis with a polarity and
  a diagnosticity value. Challenges sit on the contested claims that the
  hypotheses depend on, and contributors record attributed credences. The
  Rootclaim debate \cite{rootclaim} showed this structure, and we built
  the data model to hold it.
\item \emph{Black holes} stresses structure. A settled conclusion rests on
  a dependency chain (\code{depends\_on} relations). Each link in the
  chain stays open to inspection. The reader can see what would have to
  fail before the conclusion changes.
\item \emph{Eggs} stresses scouting. The question is open-ended. The main
  work is to find the important sub-questions. Avenues, pressure tests,
  and open cruxes carry that work.
\end{itemize}

We include the investigations, their releases, and the full receipt
trails in the submission as a database snapshot. The judge runbook
(\code{README.md}, Appendix~\ref{app:runbook}) restores the snapshot
locally in a few commands. A judge can open the rooms and inspect the
provenance panel for any claim. A judge can run a delegated brief on a
new question and can connect an MCP agent to the same graph. A judge can
inspect every claim in this paper in the artifact.

\video{--:--}{The three case-study graphs: hypotheses and credences for
COVID origins, the dependency chain for black holes, and the crux map for
eggs.}

\section{Limitations}
\label{sec:limits}

\begin{itemize}
\item \emph{Quote verification proves faithful extraction, not truth.} A
  verbatim quote from an unreliable source is still unreliable.
  Challenges, source guarantees (for example, peer-review flags), and
  attributed credences address reliability. The system verifies the chain
  of custody. The system does not claim that the chain settles the
  question.
\item \emph{Ingestion has gaps.} Paywalled sources and PDF files fall
  back to a metadata-only path. The system cannot quote-verify claims
  from that path against stored full text. The receipt records the weaker
  guarantee. The coverage gap is real for questions that depend on
  academic sources.
\item \emph{Deduplication uses a threshold.} The system decides claim
  identity with embedding similarity at a fixed cutoff. The cutoff
  sometimes merges claims that should stay distinct. The append-only rule
  limits the damage. A wrong merge adds a mention and destroys nothing,
  and \code{refines} edges restore lost distinctions. The threshold
  remains a coarse test.
\end{itemize}

\section{Related work}
\label{sec:related}

\emph{Deep research products} (OpenAI, Anthropic, Perplexity, and others)
share the fanout-read-synthesize skeleton. They do not share the output
data model. In those products, claims are not addressable, quotes are not
machine-verified, missing evidence is not an output, and two runs do not
compound. \emph{Argument mappers}, such as Kialo, have the graph, but
people curate the graph by hand. Their nodes are opinions, not
quote-receipted extractions from sources. \emph{Wikipedia} compounds and
cites sources, but it resolves disagreement into a single text. Its unit
is the article, and a reader cannot diff, merge, or credence-weight an
article. \emph{Community Notes} \cite{communitynotes} demonstrated
attributed crowd assessment at scale. It applied that assessment to
single items, not to claim graphs. \emph{Structured debate} (Rootclaim
\cite{rootclaim}) produces rich, contested material with many attributed
judgments, and we built our data model to hold that material. The 2024
COVID-origins debate is a dataset that needs this data model. We propose
a combination of git mechanics, multiplayer rooms, a receipted claim
graph, and agents as accountable contributors. To our knowledge, no other
working system combines these parts.
\section{Map to the judging dimensions}

\begin{itemize}
\item \emph{Epistemic uplift:} The system verifies each claim against a
  verbatim quote. Each node shows the evidence that supports it. The
  harness shows cruxes and thin avenues explicitly
  (\S\ref{sec:ingest}--\ref{sec:accounting}). The commons records
  attributed credences instead of one selected answer (\S\ref{sec:commons}).
\item \emph{Generalizability:} One methodology serves all cases. The
  system needs no engineering for each case. The methodology runs on
  three cases with different shapes (\S\ref{sec:demo}). Avenue planning
  is the one step that adapts to the shape of a case. This step is a
  model call, not a hand-built ontology.
\item \emph{Compounding and shareability:} The system deduplicates
  content with content addresses. The system provides commons search and
  seeding, fork, merge, release, JSON export, and a read-only MCP server
  (\S\ref{sec:compounding}, \S\ref{sec:merge}--\ref{sec:releases}). The
  outputs are graphs with receipts, not narrative summaries.
\item \emph{Scalability:} Every judgment call is a model call behind a
  schema. The budgets are constants (Table~\ref{tab:quotas}). Runs
  execute in parallel and write into one deduplicated graph. The process
  contains no hand-designed human step (\S\ref{sec:harness}).
\item \emph{Methodological transparency:} The methodology ships as
  readable skills and a state machine. Every write names its
  method@version. The methodology is therefore auditable, versioned data
  (\S\ref{sec:harness}).
\item \emph{Adversarial robustness:} The write path checks each quote in
  code and rejects invented quotes. The harness runs adversarial tests on
  every claim as a required step. Each challenge has a type and a
  receipt. No contributor can delete a challenge. \S\ref{sec:limits}
  states the failure modes.
\item \emph{Insight contribution:} The insight is the reframe itself.
  Research output becomes a multiplayer commons with git mechanics.
  Agents become accountable contributors, not tools (\S\ref{sec:platform}).
\end{itemize}

\section{Conclusion}

Epistack changes the unit of research output. The old unit is a document
that one intelligence produces. The new unit is a commons that human and
artificial intelligences maintain together.

Part 1 shows how the system makes the write path trustworthy. The harness
enforces coverage. The code checks each quote. The harness runs
adversarial tests on each claim. The system attaches a receipt to every
write.

Part 2 shows how the system makes the commons social and cumulative.
Contributors fork, review, merge, and release. Humans and agents work in
one room, with one identity model and one write path.

The system scales along stated parameters. More compute makes runs wider
and deeper. Better models improve every judgment call behind a schema.
More contributors, human or agent, deepen one deduplicated graph. The
contributors do not scatter documents.

The system runs. The three case studies are in the system. The submission
includes the snapshot. A judge can open any claim and follow its receipts.
\begin{thebibliography}{9}\small
\bibitem{flf} Future of Life Foundation.
\emph{Lab leaks, black holes, and eggs: epistemic case study competition.}
EA Forum, 2026.
\bibitem{rootclaim} P.~Miller and S.~Wilf.
\emph{The Rootclaim COVID-19 origins debate}, structured debate with judge
decisions and Bayesian analyses, 2024.
\bibitem{mcp} Anthropic. \emph{Model Context Protocol.}
\url{https://modelcontextprotocol.io}, 2024.
\bibitem{openalex} J.~Priem, H.~Piwowar, R.~Orr.
\emph{OpenAlex: a fully-open index of scholarly works.} arXiv:2205.01833,
2022.
\bibitem{communitynotes} S.~Wojcik et al.
\emph{Birdwatch: crowd wisdom and bridging algorithms can inform
understanding and reduce the spread of misinformation.} arXiv:2210.15723,
2022.
\end{thebibliography}

\appendix
\onecolumn

\section{Agent MCP tool reference}
\label{app:mcp}

The endpoint is \code{POST /api/mcp/agent/mcp} with \code{Authorization: Bearer
esk\_...} (streamable-HTTP MCP). A person mints keys in the app (``Connect an
agent''). Each key is a distinct contributor of kind \code{agent}. The system
records the person who minted the key as \code{on\_behalf\_of}. Two servers
accept read-only requests without authentication: \code{POST /api/mcp} serves
the whole commons, and \code{POST /api/mcp/<topic>/mcp} serves curated slices.

\subsection*{Read}
\begin{lstlisting}
search(query)            -> { results: [{ id, title, url }] }
fetch(id)                -> node as a document: text, verbatim quotes,
                            challenge threads, chain-of-custody receipt
get_claim(id)            -> text, modality, descriptors, evidence quotes,
                            typed relations, challenges, disputeState, receipt
get_hypothesis(id)       -> statement, linked claims (polarity, diagnosticity),
                            communityCredence, challenges, receipt
\end{lstlisting}

\subsection*{Investigations}
\begin{lstlisting}
list_investigations()    -> all rooms, newest first, with URLs and lineage
create_investigation({ title, seed_from_commons? })
                         -> new room owned by the agent; appears in every
                            member's sidebar live
get_investigation_graph({ investigation_id })
                         -> counts, hypotheses + community credence,
                            openCruxes, node catalog
\end{lstlisting}

\subsection*{Write (eve's own write path: same dedup, same receipts)}
\begin{lstlisting}
record_source({ investigation_id, text, url?, title?, author?,
                publisher?, date? })                    -> { source_id }
record_claim({ investigation_id, claim, source_id, quote,
               descriptors? })   -> { claim_id, is_new, merged_similarity }
                                    # quote MUST appear verbatim in the
                                    # source's stored text (checked in code)
record_relation({ investigation_id, from_claim_id, to_claim_id,
                  type: supports|contradicts|depends_on|refines,
                  rationale? })
record_hypothesis({ investigation_id, statement, answer_bearing? })
link_claim_to_hypothesis({ investigation_id, claim_id, hypothesis_id,
                           polarity: supports|undermines,
                           diagnosticity?: 0..1 })
record_crux({ investigation_id, claim_id, question, implication? })
record_credence({ investigation_id, hypothesis_id, credence: 0..100,
                  rationale? })     # append-only; history = belief timeline
file_challenge({ investigation_id, node_id,
                 challenge_type: counter_evidence|rival_interpretation|
                                 methodological_objection,
                 body, evidence_url? })
\end{lstlisting}

\subsection*{Collaborate}
\begin{lstlisting}
get_transcript({ investigation_id })     -> per-turn Q/A with the message ids
                                            comment anchors use
add_comment({ investigation_id, body, message_id, quote,
              quote_prefix?, quote_suffix? })   # anchored to a quote span
reply_to_comment({ comment_id, body })
send_message({ investigation_id, message, wait_for_reply? })
    # takes a REAL eve turn under the agent's name; room members watch it
    # stream exactly as they would a human member's turn
delegate_investigation({ investigation_id, brief })   -> first narration beats
delegate_step({ investigation_id, delegation_id })    -> advance one phase;
                                                         done=true at the end
\end{lstlisting}

Every write fires an activity broadcast. Clients derive agent presence and
a live agent cursor from the recency of these broadcasts. An idle agent
expires after $\sim$75\,s. Every write creates a \code{contributions} row.
The row records the contributor, the method@version, the payload hash, and
the session and turn. The row has the same shape for agent writes, human
writes, and eve writes.

\section{Worked example: one claim's chain of custody}
\label{app:receipt}

This appendix shows the provenance panel of one claim. The receipts query
resolves the panel. The shape comes verbatim from the system. We shortened
the content:

\begin{lstlisting}
claim: "..." (canonical id = hash of normalized text)
  created by  eve (agent) . method record_claim@1 . <timestamp>
              during a turn asked by <human contributor>
              in investigation "<title>"
  payload hash  sha256:...          (optional signature over the hash)
  evidence
    mention #1  source <id> "<source title>" (content-addressed)
                quote: "<verbatim passage, code-checked substring of the
                        source's stored full text>"
                retrieval: { operator: delegated_read@1, query: "...",
                             avenue: "...", via: tavily-extract,
                             delegator: <human> }
    mention #2  source <id> ...    # same claim found again elsewhere:
                                   # "many sources, one claim"
  relations   <- supports  claim <id>   "Answers the crux question: ..."
  challenges  1 open (rival_interpretation) -> derived state: contested
  credences   (on the hypothesis it supports) per-assessor timelines
\end{lstlisting}

A correction never changes this record. A superseding contribution points
to the contribution that it corrects. Readers see both contributions.

\section{Judge runbook}
\label{app:runbook}

The submission repository contains the full instructions in
\code{README.md}. The repository is a public fork with the secrets
removed. The summary follows:

\begin{enumerate}
\item Install the prerequisites: Bun, Node $\geq 24$, Docker (for local
  Supabase), and the Supabase CLI.
\item Run \code{bun install}, then \code{supabase start}. Restore the
  included case-study snapshot with one \code{psql} command. The README
  gives the exact line. Copy \code{.env.example} to \code{.env}. The file
  needs two keys: a model API key and a search API key. Run
  \code{bun run dev}.
\item Open the three case-study investigations. Inspect the provenance
  panel of any node. Scrub the belief timeline. Open a release page and
  its \code{/export} URL.
\item Delegate a brief on a question of your own. Watch the run write
  into the graph.
\item Mint an agent key in the app. Connect any MCP client to write to
  the same graph. \S\ref{sec:agents} gives the config block.
\end{enumerate}

\section{Demo video index}
\label{app:video}

One walkthrough video accompanies the submission. The
$\blacktriangleright$ boxes in the text mark the relevant timestamps.
\emph{(The submission form and the README give the link and the
timestamps.)}

\end{document}
